\section{Theory: The Conditioning Boundary}
\label{sec:theory}

The gradient-responsibility identity is local: it is a statement about the
coordinates $d$ at an exponentiation-and-normalization site. The question for a
neural network is what happens after this responsibility vector is pulled back
to parameters. This section isolates the condition under which the local
curvature bound survives that pullback, and the condition under which it is
forfeited.

\subsection{The Simplex Invariant}

\begin{lemma}[Simplex invariant and curvature bound]
\label{lem:simplex}
For
\[
    L_{\mathrm{LSE}}(d) = -\log \sum_{j=1}^K \exp(-d_j),
    \qquad
    r = \operatorname{softmax}(-d),
\]
the gradient $\nabla_d L_{\mathrm{LSE}}=r$ lies on the probability simplex.
The Hessian is
\[
    \nabla_d^2 L_{\mathrm{LSE}}
    = rr^\top - \operatorname{diag}(r).
\]
Equivalently, the negative Hessian is the categorical covariance matrix
\[
    C(r)
    = \operatorname{diag}(r)-rr^\top
    \succeq 0,
    \qquad
    \|C(r)\|_2 = \|\nabla_d^2 L_{\mathrm{LSE}}\|_2 \le \frac{1}{2}.
\]
\end{lemma}

\begin{proof}
The gradient calculation gives
\[
    \frac{\partial L_{\mathrm{LSE}}}{\partial d_j}
    =
    \frac{\exp(-d_j)}{\sum_k \exp(-d_k)}
    = r_j.
\]
Differentiating $r$ with respect to $d$ gives
\[
    \frac{\partial r_j}{\partial d_k}
    = -r_j(\mathbf{1}\{j=k\}-r_k),
\]
so $\nabla_d^2 L_{\mathrm{LSE}}=rr^\top-\operatorname{diag}(r)$. For any vector
$v$,
\[
    v^\top C(r)v
    =
    \sum_j r_j v_j^2 - \left(\sum_j r_j v_j\right)^2
    =
    \operatorname{Var}_{j\sim r}(v_j).
\]
By Popoviciu's inequality,
\[
    \operatorname{Var}_{j\sim r}(v_j)
    \le \frac{1}{4}(\max_j v_j-\min_j v_j)^2.
\]
For $\|v\|_2=1$, the squared range is at most $2$, hence
$v^\top C(r)v\le 1/2$. Taking the supremum over unit vectors proves the bound.
\end{proof}

The important feature is not merely boundedness, but \emph{where the bound
lives}. In distance coordinates it is absolute: no parameter value, scale, or
training time appears. The gradient is a normalized conditional expectation,
and the curvature is the covariance of that conditional distribution.

\subsection{Private Parameters Inherit the Bound}

\begin{proposition}[Conditioning under private affine energies]
\label{prop:private-affine}
Let each component energy be a private affine map
\[
    d_j = w_j^\top x + b_j,
\]
and let $\theta=(w_1,b_1,\ldots,w_K,b_K)$. For a single example,
\[
    \|\nabla_\theta L_{\mathrm{LSE}}\|_2
    \le \sqrt{\|x\|^2+1},
\]
and
\[
    \|\nabla_\theta^2 L_{\mathrm{LSE}}\|_2
    \le \frac{1}{2}(\|x\|^2+1).
\]
If biases are omitted from the norm or absorbed into an augmented input, this is
the bound $\|\nabla^2_\theta L_{\mathrm{LSE}}\|_2\le
\frac{1}{2}\|\tilde{x}\|^2$. Under batch-mean training,
\[
    \|\nabla_\theta^2 L_{\mathrm{LSE}}\|_2
    \le
    \frac{1}{2}\mathbb{E}\|\tilde{x}\|^2,
\]
a constant determined by the data and independent of the current parameters.
\end{proposition}

\begin{proof}
Let $\tilde{x}=(x,1)$ absorb the bias and write $d_j=\theta_j^\top\tilde{x}$.
Since $d$ is linear in the private parameter blocks, there is no second-order
term from the map $\theta\mapsto d$. The gradient block for component $j$ is
$r_j\tilde{x}$, so
\[
    \|\nabla_\theta L_{\mathrm{LSE}}\|_2^2
    =
    \sum_j r_j^2 \|\tilde{x}\|^2
    \le
    \|\tilde{x}\|^2.
\]
The Hessian has blocks
\[
    \left(\nabla_\theta^2 L_{\mathrm{LSE}}\right)_{jk}
    =
    \left(rr^\top-\operatorname{diag}(r)\right)_{jk}
    \tilde{x}\tilde{x}^\top,
\]
or, up to sign, $C(r)\otimes \tilde{x}\tilde{x}^\top$. The spectral norm of a
Kronecker product is the product of spectral norms, and Lemma~\ref{lem:simplex}
gives the factor $1/2$.
\end{proof}

Proposition~\ref{prop:private-affine} is the formal version of the
single-site conditioning claim. The smoothness constant is fixed before
training starts. Whatever step-size range is allowed by the data-dependent
constant does not shrink as the parameters move, and there is no learned
anisotropy for an adaptive optimizer to correct. This is the mechanism behind
the optimizer signature observed in the single-layer implicit-EM model.

The constant is also checkable. For a batch-mean objective with smoothness
constant $L \le \frac{1}{2}\mathbb{E}\|\tilde{x}\|^2$, the descent lemma gives
a stable full-batch gradient-descent scale on the order of $2/L$. The
experiments therefore report $\mathbb{E}\|\tilde{x}\|^2$ and compare this scale
to the learning-rate tolerance observed in the stop-gradient sweep.

\begin{proposition}[Smooth monotone kernels]
\label{prop:smooth-kernel}
Let
\[
    d_j = \phi(w_j^\top x+b_j)
\]
with private component parameters, $\phi'\ge 0$, $|\phi'|\le B_1$, and
$|\phi''|\le B_2$. Then the parameter-space curvature of
$L_{\mathrm{LSE}}(d(\theta))$ is bounded by a constant multiple of
$\|\tilde{x}\|^2$ that depends only on $B_1$ and $B_2$, not on the current
parameters.
\end{proposition}

\begin{proof}
The chain rule gives two terms:
\[
    \nabla_\theta^2 L
    =
    J^\top \nabla_d^2 L J
    +
    \sum_j \frac{\partial L}{\partial d_j}\nabla_\theta^2 d_j.
\]
The first term is bounded by Lemma~\ref{lem:simplex} and
$\|J\|^2\le B_1^2\|\tilde{x}\|^2$. The second term is block diagonal across
components; since $\partial L/\partial d_j=r_j$ and $\sum_j r_j=1$, its norm is
bounded by $B_2\|\tilde{x}\|^2$. Thus the total curvature is bounded by
$(B_1^2/2+B_2)\|\tilde{x}\|^2$ up to the harmless convention for including
biases.
\end{proof}

Softplus is the main smooth case in this paper. ReLU should be read separately:
its curvature bounds hold almost everywhere, but it introduces a zeroth-order
absorbing-state pathology. When a unit is gated off, its gradient can become
exactly zero, and the unit may not recover. That failure is not a violation of
the curvature bound; it is the neural version of hard-assignment cluster death.

\subsection{Composition Forfeits the Uniform Bound}

We now compose the same local site behind a learned dense map. Let
\[
    y = \operatorname{NLS}(d),
    \qquad
    h = W_2y,
    \qquad
    L_{\mathrm{CE}} = -h_c + \log\sum_{\ell}\exp(h_\ell),
\]
and $p=\operatorname{softmax}(h)$. The Jacobian of
$y=\operatorname{NLS}(d)$ is
\[
    J = \frac{\partial y}{\partial d}
      = I - \mathbf{1}r^\top,
\]
because $\partial \log \sum_k \exp(-d_k)/\partial d_j=-r_j$.

\begin{proposition}[Composition forfeits the uniform bound]
\label{prop:composition}
For the composed supervised path above, the Hessian of $L_{\mathrm{CE}}$ with
respect to the intermediate distances $d$ contains the conjugated softmax
Hessian
\[
    J^\top W_2^\top
    \left(\operatorname{diag}(p)-pp^\top\right)
    W_2J,
\]
plus curvature terms from the NLS map. Consequently, its norm is controlled by
a bound whose leading term scales as $\sigma_{\max}(W_2)^2$ rather than by a
parameter-independent constant.
\end{proposition}

\begin{proof}
The CE Hessian with respect to logits is
\[
    H_h L_{\mathrm{CE}}
    =
    \operatorname{diag}(p)-pp^\top.
\]
Pulling this term back through the linear map $h=W_2y$ and the NLS Jacobian
$J$ gives the displayed conjugated term. The remaining terms come from the
second derivatives of $y=\operatorname{NLS}(d)$, weighted by the incoming
gradient from the head. Since
\[
    \|W_2^\top(\operatorname{diag}(p)-pp^\top)W_2\|_2
    \le
    \frac{1}{2}\sigma_{\max}(W_2)^2,
\]
up to the NLS Jacobian constants, the bound now depends on the learned spectrum
of $W_2$.
\end{proof}

Proposition~\ref{prop:composition} is deliberately a
guarantee-forfeiture statement. It shows that the parameter-independent
$1/2$ bound of the local LSE site no longer controls the composed supervised
path. It does \emph{not} prove that the curvature must be large for every
possible $W_2$ or every point in training. The empirical curvature measurements
in Section~\ref{sec:mechanism} carry that claim for the studied networks: the
composed CE-path curvature is larger, non-monotone, and parameter-dependent.

We state the clean proposition without LayerNorm. The measurements are taken on
the actual model, including LayerNorm, so its effect is included in the reported
curvature. A proof with LayerNorm would insert another per-sample Jacobian and
additional input-dependent terms; it would not restore the parameter-independent
LSE guarantee.

\paragraph{Scope of the claim.}
The propositions above establish sufficient conditions for preserving the local
EM conditioning guarantee: private component parameters, an unmixed LSE
objective, and monotone per-coordinate kernels. They also show that an
intervening shared, sign-indefinite learned map forfeits the uniform guarantee.
They do not prove necessity for arbitrary architectures and do not provide a
convergence theorem for the composed system.

\subsection{Operations That Preserve or Destroy the Invariant}

The preceding propositions can be summarized by how operations treat the
simplex-valued responsibility gradient.

\paragraph{Preserving operations.}
Private per-component parameters preserve the alignment between responsibility
$r_j$ and parameter block $\theta_j$. Per-coordinate monotone maps preserve sign
and component identity, changing only local scales. Smooth bounded kernels
preserve parameter-independent curvature up to constants. Stochastic maps are a
partial first-order preservation class: if $A$ is row-stochastic, then pulling a
nonnegative unit-mass vector back through $A^\top$ preserves nonnegativity and
total mass because $\mathbf{1}^\top A^\top g=(A\mathbf{1})^\top g=\mathbf{1}^\top
g$.

\paragraph{Destroying operations.}
Dense sign-indefinite learned maps do not preserve nonnegativity, normalization,
or component alignment. Objective mixing replaces a single posterior gradient
with a sum of gradients from different objectives; the sum is generally not the
posterior of any one latent-variable model. Hard gates such as ReLU can turn
small responsibility into exactly zero gradient, creating absorbing states.
Shared lower parameters couple the component updates and destroy the separated
M-step structure.

\subsection{Classical EM Connection}

The last point is the link to classical EM. In the clean case, the E-step
produces responsibilities and the expected complete-data objective separates:
\[
    Q(\theta)=\sum_{j=1}^K Q_j(\theta_j).
\]
Each component owns its own parameters, and the M-step decouples into
independent per-component optimizations. This separation is the source of EM's
algorithmic structure.

If all components depend on a shared lower map, the separation disappears.
Writing the lower parameters as $\psi$, the complete-data terms have the form
$Q_j(\theta_j,\psi)$, so
\[
    Q(\theta,\psi)=\sum_j Q_j(\theta_j,\psi)
\]
is no longer separable in $\psi$. The M-step becomes a single coupled
optimization problem. Backpropagation through shared lower layers is the neural
version of this classical shared-parameter failure, and the conjugated Hessian
in Proposition~\ref{prop:composition} is its coordinate-space fingerprint.

\begin{table}[h]
\centering
\begin{tabular}{p{0.25\linewidth}p{0.32\linewidth}p{0.32\linewidth}}
\toprule
Classical EM failure & Why EM weakens & Neural incarnation \\
\midrule
Shared component parameters & $Q$ stops separating; the M-step is coupled &
Backpropagation through shared lower maps \\
Non-separable / generalized EM & No closed-form independent M-step &
Component energies computed by deep maps \\
Singular components & Variance collapses and likelihood becomes unbounded &
LSE-only collapse without volume control \\
Hard assignment & Zero-responsibility components cannot recover &
Dead ReLU units \\
High missing information & Diffuse responsibilities slow convergence &
Low $1/2$-saturation and weak competition \\
\bottomrule
\end{tabular}
\caption{Classical EM failure modes and their neural counterparts.}
\label{tab:em-failures}
\end{table}

This structural story must be reconciled with the graded empirical result. The
claim is not that sharing alone always destroys the effective conditioning of
the full update. Rather, sharing creates a foreign CE-path contribution whose
curvature is no longer controlled by the local LSE bound. The update at $W_1$
contains both this contribution and the local EM contribution:
\[
    \nabla_{W_1} L_{\mathrm{total}}
    =
    \nabla_{W_1} L_{\mathrm{CE}}
    +
    \lambda \nabla_{W_1}L_{\mathrm{aux}}.
\]
When the CE contribution dominates, the effective landscape is the
non-separable shared-parameter landscape. When the EM contribution dominates,
the same graph can behave like the isolated EM site. Thus the structural result
explains why the foreign term is ill-conditioned, while the dominance result
explains when that term matters.
